\documentclass[10pt,a4paper]{article}
\usepackage{amsmath,amssymb,amsthm,graphicx}
\usepackage{hyperref}
% Card-local minimal preamble (journal unit sigma-2023-014-a-generalization-of-zwegers-mu-function-acco, P2-ahl16 precedent).
%% Verbatim from the pinned source: the author theorem environments and macro block,
%% so every macro and environment of the counted spans is available.
% Omitted and why: the documentclass line (the journal class sigma is not on the host;
%% the assembler supplies the class) and page-geometry/font blocks (they fight the
%% forced xelatex engine).
\usepackage{amsmath,amssymb,amsthm,mathtools}

\numberwithin{equation}{section}
\newtheorem{thm}{Theorem}[section]
\newtheorem{prop}[thm]{Proposition}
\newtheorem{lem}[thm]{Lemma}
\newtheorem{cor}[thm]{Corollary}
\newtheorem{fact}[thm]{Fact}
\newtheorem{conj}[thm]{Conjecture}
\newtheorem{prob}[thm]{Problem}
\newtheorem{ques}[thm]{Question}
\theoremstyle{definition}
\newtheorem{dfn}[thm]{Definition}
\newtheorem{rmk}[thm]{Remark}
\newtheorem{exa}[thm]{Example}
\newcommand{\Hz}{\mathbb{H}}
\newcommand{\Zz}{\mathbb{Z}}
\newcommand{\Cz}{\mathbb{C}}
\newcommand{\hyper}[4]{\left(\begin{matrix} #1 \\ #2 \end{matrix}; #3, #4\right)}
\newcommand{\sn}{\mathrm{sn}\,}
\newcommand{\cn}{\mathrm{cn}\,}
\newcommand{\dn}{\mathrm{dn}\,}
\title{The q-Borel plus q-Laplace image of the q-Hermite-Weber fundamental solution equals the generalized mu-function}
\author{Shibukawa, Genki, Tsuchimi, Satoshi}
\date{doi:10.3842/SIGMA.2023.014; CC BY 4.0}
\begin{document}
\maketitle
\setcounter{section}{1}
\setcounter{thm}{1}
\setcounter{equation}{21}
\noindent\textit{Editorial note.} All delivered mathematical text is the author's own, copied from the pinned journal source file sources/sigma-2023-014-a-generalization-of-zwegers-mu-function-acco/pinned-original-source.tex (A Generalization of Zwegers' mu-Function According to the q-Hermite-Weber Difference Equation); no statement or proof step was rewritten, reordered or removed.\par\medskip
\noindent\textit{Editorial note.} The counted proof is the author's complete argument: Seed convention: counters stepped from the source itself (tmp-step-2023 walk, cross-validated against the host aux ground truth of the pinned source compiled whole): section 1 of the shared [thm] counter holds one environment (def:mua 1.1) before this theorem, and section 1 holds numbered equations (1.1)-(1.21) before the span, so the card prints Theorem 1.2 with in-span equation (1.22) eq:mu and q-Hermite--Weber (label line 379). The proof is the author's own 'Proof of Theorem' environment (source lines 746-780) delivered verbatim including the optional-argument back-reference to the theorem and the tag-qed plus renewcommand-qed device; all its displays are unnumbered (align* and displayed equations), so the card carries no further numbered equations after (1.22). The statement sits in Section 1 and the proof is delivered in Section 2 of the article; since the proof carries no numbered equations, the single-seed card numbering reproduces the article's printed numbers exactly. Out-of-span navigation pointers: none (no \textbackslash{}ref outside the counted spans). Citations: the proof carries cite GR once (line 747, q-binomial theorem p. 8 equation (1.3.2)); resolved by the retained-bibliography recipe (duty ruling 2026-10-14, per Chief instruction to lane-x2, 022 precedent L48545): the entry is copied verbatim from the author's own thebibliography (source lines 1531-1533) into the card-local fragment retained-bibliography.tex with bibliography backend none.\par\medskip
\noindent\textit{Editorial note.} The label thm:mu and q-Hermite--Weber names the counted statement; its proof is the author's own 'Proof of Theorem' environment delivered in Section 2 of the article (source lines 746-780), retained verbatim as the card's proof. The statement and the proof carry no numbered equations beyond the in-span (1.22), so the single-seed card numbering reproduces the article's printed numbers exactly.\par\medskip
\noindent\textit{Editorial note.} The counted statement is carried by the author's own thm environment in the article source.\par\medskip
\noindent\textit{Editorial note.} Retained-bibliography recipe (duty ruling 2026-10-14, per Chief instruction to lane-x2; 022 precedent L48545): the literature citation carried by the delivered span is disclosed in the notes of this card and delivered verbatim as a card-local bibliography fragment: entry GR (source lines 1531-1533, cited at line 747 for the q-binomial theorem, p. 8, equation (1.3.2)), copied without any re-typing into focused-seeds/combinatorics/X2-sigma-2023-014-thm-mu-and-q-Hermite-Weber/retained-bibliography.tex, with bibliography backend none and the hyperref package for the fragment's links. The cited q-binomial theorem is a standard published input used inside the proof, per the declared-dependency class of the S12-021 published-input ruling.\par\medskip
\noindent\textit{Editorial note.} PRESCREEN RE-READ: cite\_missing resolved by the retained-bibliography recipe (duty ruling 2026-10-14, per Chief instruction to lane-x2, 022 precedent L48545): the cited entries are copied verbatim from the author's own thebibliography of the pinned unit source into the card-local fragment retained-bibliography.tex with bibliography backend none; no citation is dropped or re-typed.\par\medskip
\noindent\textit{Editorial scope.} the counted statement and its own proof environment only\par\medskip
\noindent\textit{Allowed standard prerequisites (editorial).}\begin{itemize}
\item The series definition of the generalized mu-function def:mua (1.1, source lines 343-362), the value the proof identifies at the end.
\item The fundamental solution of the q-Hermite-Weber equation as the 2phi0 series t-tilde-f\_0 (source lines 365-369) and the q-Laplace transformation eq:q-Laplace (1.21, source line 331) composing it into f\_0.
\item The q-binomial theorem, cited at line 747 from entry GR (p. 8, equation (1.3.2)), a standard published input used to sum the q-Borel image.
\end{itemize}\medskip
\begin{thm}\label{thm:mu and q-Hermite--Weber}
Let $f_{0}(x,\lambda )$ be the image of the $q$-Borel and $q$-Laplace transformations of the fundamental solution $\widetilde{f}_{0}$ of the $q$-Hermite--Weber equation at $x=0$:
\[
f_{0}(x)
 :=
 x^{\frac{\alpha }{2}}\mathcal{L}^{+}\circ\mathcal{B}^{+}\big(\widetilde{f}_{0}\big)(x,\lambda).
\]
Then we have
\begin{gather}\label{eq:mu and q-Hermite--Weber}
f_{0}\big({\rm e}^{2\pi{\rm i}(u-v)},-{\rm e}^{2\pi{\rm i}u}\big)
 =\mu(u,v;\alpha).
\end{gather}
\end{thm}

\begin{proof}[Proof of Theorem \ref{thm:mu and q-Hermite--Weber}]
From the definition of the $q$-Borel transformation and the $q$-binomial theorem (see \cite[p.~8, equation~(1.3.2)]{GR}),
\begin{align*}
\mathcal{B}^{+}\left(\widetilde{f}_{0}\right)(\xi )
 &=
 \sum_{n\geq 0}\frac{(a)_{n}}{(q)_{n}}(-1)^{n}\left(\frac{\xi }{a}\right)^{n}
 =
 \frac{(-\xi )_{\infty }}{(-\xi /a)_{\infty }}.
\end{align*}
By a simple calculation, we have
\begin{align*}
\mathcal{L}^{+}\circ\mathcal{B}^{+}\big(\widetilde{f}_{0}\big)(x,-\lambda)
 &=
 \sum_{n\in\Zz}\frac{1}{\theta_q(-\lambda q^n/x)}\frac{(\lambda q^n)_\infty}{(\lambda q^n/a)_\infty} \\
 &=
 \frac{1}{\theta_q(-\lambda/x)}\sum_{n\in\Zz}\left(-\frac{\lambda}{x}\right)^{n+1}q^\frac{n(n+1)}{2}
 \frac{\big(\lambda q^{n+1}\big)_{\infty }}{\big(\lambda q^{n-\alpha +1}\big)_{\infty }}.
\end{align*}
The second equality follows from the well-known $q$-difference relation of the theta function:
\[
\theta _{q}(q^{n}x)=q^{-\frac{n(n-1)}{2}}x^{-n}\theta _{q}(x), \qquad n \in \mathbb{Z}.
\]
By substituting $x={\rm e}^{2\pi{\rm i}(u-v)}$ and $\lambda={\rm e}^{2\pi{\rm i}u}$, we have
\begin{align*}
f_0\big({\rm e}^{2\pi{\rm i}(u-v)},-{\rm e}^{2\pi{\rm i}u}\big)
 &=
 {\rm e}^{\pi{\rm i}\alpha(u-v)}\mathcal{L}^{+}\circ\mathcal{B}^{+}\big(\widetilde{f}_{0}\big)\big({\rm e}^{2\pi{\rm i}(u-v)},-{\rm e}^{2\pi{\rm i}v}\big) \\
 &=
 -\frac{{\rm e}^{\pi{\rm i}\alpha(u-v)}}{\theta_q\big({-}{\rm e}^{2\pi{\rm i}v}\big)}\sum_{n\in\Zz}(-1)^n{\rm e}^{2\pi{\rm i}(n+1)v}q^\frac{n(n+1)}{2}\frac{\big({\rm e}^{2\pi{\rm i}u}q^{n+1}\big)_\infty}{\big({\rm e}^{2\pi{\rm i}u}q^{n-\alpha+1}\big)_\infty} \\
 &=
 \frac{{\rm i}{\rm e}^{\pi{\rm i}\alpha(u-v)}}{\vartheta_{11}(v)}\sum_{n\in\Zz}(-1)^n{\rm e}^{2\pi{\rm i}\left(n+\frac{1}{2}\right)v}q^\frac{\left(n+\frac{1}{2}\right)^2}{2}\frac{\big({\rm e}^{2\pi{\rm i}u}q^{n+1}\big)_\infty}{\big({\rm e}^{2\pi{\rm i}u}q^{n-\alpha+1}\big)_\infty} \\
 &=
 {\rm i}q^\frac{1}{8}\mu(u,v;\alpha).\tag*{\qed}
\end{align*}\renewcommand{\qed}{}
\end{proof}

%% Card-local bibliography fragment: the author's own \bibitem entry copied
%% verbatim from the pinned unit source (bibliography_backend none).
\begin{thebibliography}{99}
\bibitem{GR}
Gasper G., Rahman M., Basic hypergeometric series, 2nd ed., \textit{Encyclopedia Math.
 Appl.}, Vol.~96, \href{https://doi.org/10.1017/CBO9780511526251}{Cambridge University Press}, Cambridge, 2004.

\end{thebibliography}
\end{document}
